% الجزء: نظرية البرهان
% الفصل: حذف القطع
% القسم: ce-largest

\documentclass[../../../include/open-logic-section]{subfiles}

\begin{document}

\olfileid{pt}{cut}{inv}

\olsection{إزالة القطوع ذات الرتبة القصوى}

يبيّن برهان \olref[seq][inv]{lem:G3c-invert} أنه إذا كان النسق \Log{G3c}
!!{prove}s نتيجة قاعدة منطقية (غير
\RightR{\lexists} و\LeftR{\lforall})، فإنه !!{prove}s أيضًا كل مقدمة من
مقدمات القاعدة من دون زيادة !!{height} البرهان. ويمكننا إثبات نتيجة أقوى،
وهي أن هذا يصدق أيضًا على $\Log{G3c} + \CutCS$.

وكما سبق، نعرّف \emph{رتبة القطع} لاستدلال \CutCS{} بأنها عمق !!{formula}
القطع فيه.

\begin{lem}\ollabel{lem:inv-G3c-cut}
إذا كان $\Log{G3c} + \CutCS$ !!{prove}s نتيجة قاعدة منطقية~$R$ غير
\RightR{\lexists} أو \LeftR{\lforall} باستعمال !!a{proof}~$\pi$، فإنه
!!{prove}s أيضًا كل مقدمة بــ!!{proof}~$\pi'$ (أو بــ!!{proof}s~$\pi'$،
$\pi''$ بحسب كون القاعدة ذات مقدمة واحدة أو مقدمتين). وفضلًا عن ذلك،
(أ)~$\pheight{\pi'}\le\pheight{\pi}$ (وكذلك
$\pheight{\pi''}\le\pheight{\pi}$)، و(ب)~يكون عدد استدلالات \CutCS{}
ورتبها في~$\pi'$ (وفي $\pi''$) هو نفسه في~$\pi$.
\end{lem}

\begin{proof}
يتضح ذلك بفحص براهين
\cref{pt:seq:inv:lem:G3c-invert,pt:seq:inv:lem:invert-quant}. فقد كان ذلك
البرهان استقراء على~$\pheight{\pi}$. وفي حالة الأساس، استبدلنا مسلّمة
بمسلّمة أخرى، ومن ثم يتحقق الشرطان (أ) و(ب) بداهة. وإذا كانت القاعدة
الأخيرة في~$\pi$ هي القاعدة المطلوبة~$R$، فإن !!{proof}sها الفرعية
المباشرة هي !!{proof}s~$\pi_i$ التي نريدها، ويتحقق الشرطان (أ) و(ب). وفي
جميع الحالات الأخرى،
طبّقنا فرض الاستقراء على !!{proof}s~$\pi$ الفرعية المباشرة، أي التي تنتهي
بمقدمة أو مقدمات الاستدلال الأخير~$S$، ثم أعدنا تطبيق القاعدة نفسها~$S$ على
النتائج. ويتحقق الشرطان (أ) و(ب) في !!{proof}s البرهان الجديد الفرعية
المباشرة، ولذلك يتحققان في البرهان كله. ويبقى التحقق من الحالة التي يكون
فيها الاستدلال الأخير هو~\CutCS. ونكتفي مرة أخرى بمثال واحد: افترض أن
$\pi$ هو !!a{proof} لنتيجة \LeftR{\land}، وهي $!B \land !C, \Gamma
\Sequent \Delta$، وأنه ينتهي بــ~\CutCS:
\[
\AxiomC{}
\RightLabel{$\pi_1$}
\Deduce$!B \land !C, \Gamma \fCenter \Delta, !A$
\AxiomC{}
\RightLabel{$\pi_2$}
\Deduce$!A, !B \land !C, \Gamma \fCenter \Delta$
\RightLabel{\CutCS}
\BinaryInf$!B \land !C, \Gamma \fCenter \Delta$
\DisplayProof
\]
ينطبق فرض الاستقراء على $\pi_1$ و~$\pi_2$ (وهما أيضًا !!{proof}s لأمثلة
من نتيجة $\LeftR{\land}$)، ويعطينا !!{proof}s~$\pi_1'$ و~$\pi_2'$ لكل من
$!B, !C, \Gamma \fCenter \Delta, !A$ و$!A, !B, !C, \Gamma \fCenter
\Delta$ على الترتيب. ونجمعهما باستعمال \CutCS{} لنحصل على~$\pi'$:
\[
\AxiomC{}
\RightLabel{$\pi_1'$}
\Deduce$!B, !C, \Gamma \fCenter \Delta, !A$
\AxiomC{}
\RightLabel{$\pi_2'$}
\Deduce$!A, !B, !C, \Gamma \fCenter \Delta$
\RightLabel{\CutCS}
\BinaryInf$!B, !C, \Gamma \fCenter \Delta$
\DisplayProof
\]
ومن الواضح أن الشرطين (أ) و(ب) متحققان.
\end{proof}

لاحظ أن هذا لا يصلح إذا استعملنا \Cut{} بدلًا من \CutCS، لأن نتيجة
!!{proof} الأخير ستحتوي عندئذ نسختين من $!B, !C, \Gamma$ في المقدم ونسختين
من $\Delta$ في التالي. وسنضطر إلى الاحتجاج بمقبولية الانكماش، لكن برهان
\olref[seq][inv]{prop:G3c-cont-adm} يستعمل مبرهنة قابلية العكس، فنكون بذلك
قد استدللنا استدلالًا دائريًا.

يمكننا استعمال \cref{pt:cut:inv:lem:inv-G3c-cut} لإعطاء برهان بديل لحذف
القطع. ففي هذا البرهان لا نزيل تباعًا المواضع العليا لاستدلالات \CutCS{}،
بل نزيل مواضع استدلالات \CutCS{} ذات الرتبة القصوى. أي إننا نبدأ بــ
!!a{proof}~$\pi$ في $\Log{G3c} + \CutCS$، ونزيل قطعًا في أي موضع منه
(وقد نستبدله بقطوع أصغر رتبة)، ثم نكرر ذلك إلى أن لا يبقى أي قطع. والشرط
الوحيد هو ألا يقع فوق القطع الذي نعالجه قطع من الرتبة نفسها أو من رتبة أعلى.

\begin{lem}\ollabel{lem:max-cut-red-G3c} إذا كان $\pi$ هو !!a{proof} في
$\Log{G3c}+\CutCS$ ينتهي باستدلال~\CutCS{} رتبته~$n$، ولا يستعمل فيما عدا
ذلك إلا قواعد~$\Log{G3c}$ واستدلالات \CutCS{} رتبتها $<n$، فهناك برهان
$\pi'$ للمتتالية النهائية نفسها في~$\Log{G3c} + \CutCS$ تكون رتبة كل
استدلال \CutCS{} فيه $<n$.
\end{lem}

\begin{proof}
ينتهي !!{proof}~$\pi$ بما يأتي:
\[
\AxiomC{}
\RightLabel{$\pi_1$}
\Deduce$\Gamma \fCenter \Delta, !A$
\AxiomC{}
\RightLabel{$\pi_2$}
\Deduce$!A, \Gamma \fCenter \Delta$
\RightLabel{\CutCS}
\BinaryInf$\Gamma \fCenter \Delta$
\DisplayProof
\]
نميّز حالات بحسب صورة~$!A$.

افترض أن $!A$ ذرية. لاحظ أنه، بمقتضى الشرط القائل إن كل استدلال \CutCS{}
في~$\pi_1$ و$\pi_2$ لا بد أن تكون رتبته $< \depth{!A} = 0$، لا يمكن أن يوجد أي
استدلال \CutCS{} في~$\pi_1$ أو~$\pi_2$. نبيّن بالاستقراء على
$\pheight{\pi_2}$ أن لـ$\Gamma \Sequent \Delta$ برهانًا خاليًا من القطع.
إذا كان $\pheight{\pi_2}=0$، فإن $!A, \Gamma \Sequent \Delta$ مسلّمة.
فإذا لم تكن $!A$ رئيسة، كان بعض $!C$ الذري !!a{element} في كل من~$\Gamma$
و$\Delta$، وكانت $\Gamma \Sequent \Delta$ نفسها مسلّمة. وإذا كانت $!A$
رئيسة، فإن $\Delta=\Delta', !A$. وفي هذه الحالة ينتهي $\pi_1$ بــ
$\Gamma \Sequent \Delta', !A, !A$. ونحصل على برهان خال من القطع لــ
$\Gamma \Sequent \Delta', !A$ بتطبيق
\olref[seq][inv]{prop:G3c-cont-adm} (مقبولية الانكماش). وإذا كان
$\pheight{\pi_2}>0$، فإنه ينتهي بقاعدة~$S$. خذ مقدمة من مقدمات تلك القاعدة؛
صورتها $!A, \Gamma', \Pi \Sequent \Delta', \Lambda$، حيث $\Pi$ و$\Lambda$
هما !!{formula}s الفعالة للقاعدة~$S$. يعطي تطبيق
\olref[seq][adm]{prop:weak-G3c-adm} على~$\pi_1$ و$\pi_2$ !!{proof}s لكل من
$\Gamma, \Gamma', \Pi \Sequent \Delta, \Delta', \Lambda, !A$ و
$!A, \Gamma, \Gamma', \Pi \Sequent \Delta, \Delta', \Lambda$، وينطبق
عليهما فرض الاستقراء. وهكذا نحصل على !!a{proof} لــ$\Gamma, \Gamma', \Pi
\Sequent \Delta, \Delta', \Lambda$، لكل مقدمة من مقدمات~$S$. وبتطبيق $S$
نحصل على $\Gamma, \Gamma \Sequent \Delta, \Delta$، ثم يزوّدنا
\olref[seq][inv]{prop:G3c-cont-adm} بالبرهان الخالي من القطع لــ$\Gamma
\Sequent \Delta$.

إذا لم تكن $!A$ ذرية، نميّز حالات بحسب !!{main operator} لــ$!A$. وفي كل
حالة نطبّق \olref{lem:inv-G3c-cut} لنحصل على !!{proof}s لمقدمات قاعدتي
اليسار واليمين اللتين تكون $!A$ صيغتهما الرئيسة. ثم نتابع كما في الحالة~E
من برهان \olref[top]{lem:cut-adm-G3c}.
\begin{enumerate}
  \item $!A \ident \lnot !B$. ينتهي !!{proof}s~$\pi_1$ و$\pi_2$ بكل من
  $\Gamma \Sequent \Delta, \lnot !B$ و$\lnot !B, \Gamma \Sequent
  \Delta$. وبحسب \olref{lem:inv-G3c-cut} يوجد !!{proof}s~$\pi_1'$ و
  $\pi_2'$ لــ$!B, \Gamma \Sequent \Delta$ و$\Gamma \Sequent \Delta,
  !B$. ويكون !!{proof}~$\pi'$ كما يأتي:
  \[
  \AxiomC{}
  \RightLabel{$\pi_2'$}
  \Deduce$\Gamma \fCenter \Delta, !B$
  \AxiomC{}
  \RightLabel{$\pi_1'$}
  \Deduce$!B, \Gamma \fCenter \Delta$
  \RightLabel{\CutCS}
  \BinaryInf$\Gamma \fCenter \Delta$
  \DisplayProof
  \]
  \item $!A \ident (!B \lor !C)$. ينتهي !!{proof}s~$\pi_1$ و$\pi_2$
  بكل من $\Gamma \Sequent \Delta, !B \lor !C$ و$!B \lor !C, \Gamma
  \Sequent \Delta$. وبحسب \olref{lem:inv-G3c-cut} (قابلية عكس
  \RightR{\lor} مطبقة على~$\pi_1$)، يوجد !!a{proof}~$\pi_1'$ لــ
  $\Gamma \Sequent \Delta, !B, !C$. وبحسب \olref{lem:inv-G3c-cut}
  (قابلية عكس \LeftR{\lor} مطبقة على~$\pi_2$)، يوجد !!a{proof}~$\pi_2'$
  لــ$!B, \Gamma \Sequent \Delta$ و!!a{proof}~$\pi_2''$ لــ$!C, \Gamma
  \Sequent \Delta$. وبتطبيق \olref[seq][adm]{prop:weak-G3c-adm} على
  $\pi_2''$ نحصل أيضًا على !!a{proof}~$\pi_2'''$ لــ$!C, \Gamma \Sequent
  \Delta, !B$. ويكون !!{proof}~$\pi'$:
  \[
  \AxiomC{}
  \RightLabel{$\pi_1'$}
  \Deduce$\Gamma \fCenter \Delta, !B, !C$
  \AxiomC{}
  \RightLabel{$\pi_2'''$}
  \Deduce$!C, \Gamma \fCenter \Delta, !B$
  \RightLabel{\CutCS}
  \BinaryInf$\Gamma \fCenter \Delta, !B$
  \AxiomC{}
  \RightLabel{$\pi_2'$}
  \Deduce$!B, \Gamma \fCenter \Delta$
  \RightLabel{\CutCS}
  \BinaryInf$\Gamma \fCenter \Delta$
  \DisplayProof
  \]
  \item $!A \ident (!B \land !C)$. تمرين.
  \item $!A \ident (!B \lif !C)$. ينتهي !!{proof}s~$\pi_1$ و$\pi_2$
  بكل من $\Gamma \Sequent \Delta, !B \lif !C$ و$!B \lif !C, \Gamma
  \Sequent \Delta$. وبحسب \olref{lem:inv-G3c-cut} (قابلية عكس
  \RightR{\lif} مطبقة على~$\pi_1$)، يوجد !!a{proof}~$\pi_1'$ لــ
  $!B, \Gamma \Sequent \Delta, !C$. وبحسب \olref{lem:inv-G3c-cut}
  (قابلية عكس \LeftR{\lif} مطبقة على~$\pi_2$)، يوجد !!a{proof}~$\pi_2'$
  لــ$\Gamma \Sequent \Delta, !B$ و!!a{proof}~$\pi_2''$ لــ$!C, \Gamma
  \Sequent \Delta$. وبتطبيق \olref[seq][adm]{prop:weak-G3c-adm} على
  $\pi_2''$ نحصل أيضًا على !!a{proof}~$\pi_2'''$ لــ$!C, !B, \Gamma
  \Sequent \Delta$. ويكون !!{proof}~$\pi'$:
  \[
  \AxiomC{}
  \RightLabel{$\pi_2'$}
  \Deduce$\Gamma \fCenter \Delta, !B$
  \AxiomC{}
  \RightLabel{$\pi_1'$}
  \Deduce$!B, \Gamma \fCenter \Delta, !C$
  \AxiomC{}
  \RightLabel{$\pi_2'''$}
  \Deduce$!C, !B, \Gamma \fCenter \Delta$
  \RightLabel{\CutCS}
  \BinaryInf$!B, \Gamma \fCenter \Delta$
  \RightLabel{\CutCS}
  \BinaryInf$\Gamma \fCenter \Delta$
  \DisplayProof
  \]
  \item $!A \ident \lforall[x][!B(x)]$. ينتهي !!{proof}s~$\pi_1$ و
  $\pi_2$ بكل من $\Gamma \Sequent \Delta, \lforall[x][!B(x)]$ و
  $\lforall[x][!B(x)], \Gamma \Sequent \Delta$. وبحسب
  \olref{lem:inv-G3c-cut} يوجد !!a{proof}~$\pi_1'$ لــ$\Gamma \Sequent
  \Delta, !B(c)$.

  تأمل الآن !!{proof}~$\pi_2$. إنه ينتهي بــ$\lforall[x][!B(x)],
  \Gamma \Sequent \Delta$. وبينما نتحرك صعودًا ابتداء من المتتالية النهائية
  لــ$\pi_2$، نعلّم كل ورود لــ$\lforall[x][!B(x)]$ على يسار متتالية «يقود
  إلى» ورود $\lforall[x][!B(x)]$ في المتتالية النهائية لــ$\pi_2$، أي:
  (أ)~علّم $\lforall[x][!B(x)]$ في المتتالية النهائية لــ$\pi_2$.
  (ب)~إذا وردت $\lforall[x][!B(x)]$ في سياق نتيجة قاعدة وكانت معلّمة،
  فعلّم الورود المناظر في مقدمة القاعدة أو مقدماتها. (ج)~إذا كانت
  $\lforall[x][!B(x)]$ رئيسة في نتيجة قاعدة \LeftR{\lforall} وكانت معلّمة،
  فعلّم الورودين الفعالين المناظرين لــ$\lforall[x][!B(x)]$ و$!B(t)$ في
  المقدمة.

  تأمل الآن كل هذه الورودات المعلّمة لــ$\lforall[x][!B(x)]$. لا يكون أي
  منها فعالًا في قاعدة غير \LeftR{\lforall} (فلا تُعلّم إلا !!{formula}s
  الفعالة لقاعدة \LeftR{\lforall}). احذف من~$\pi_2$ كل الورودات المعلّمة
  لــ$\lforall[x][!B(x)]$.
  إذا كان الناتج !!{proof} صحيحًا، فلم تكن الورودات المعلّمة لــ
  $\lforall[x][!B(x)]$ فعالة قط، بل أتت كلها مباشرة من مسلّمات. وينتهي
  !!{proof} بــ$\Gamma \Sequent \Delta$، ويمكننا أن نستبدل به $\pi$.
  وبذلك نكون قد أزلنا قطعًا ذا رتبة قصوى.

  وإلا، فالناتج ليس برهانًا صحيحًا، لأننا حذفنا !!{formula}s
  $\lforall[x][!B(x)]$ كانت فعالة في بعض استدلالات \LeftR{\lforall}.
  وقد حوّلنا هذه إلى «استدلالات» من الصورة
  \[
  \AxiomC{}
  \RightLabel{$\pi_t$}
  \Deduce$!B(t), \Pi \fCenter \Lambda$
  \RightLabel{\LeftR{\lforall}}
  \UnaryInf$\Pi \fCenter \Lambda$
  \DisplayProof
  \]
  استبدل بكل «استدلال» علوي من هذا النوع ما يأتي:
  \[
  \AxiomC{}
  \RightLabel{$\Subst{\pi_1'}{t}{c}[\Pi \Sequent \Lambda]$}
  \Deduce$\Gamma, \Pi \fCenter \Delta, \Lambda, !B(t)$
  \AxiomC{}
  \RightLabel{$\pi_t[\Gamma \Sequent \Delta]$}
  \Deduce$!B(t), \Gamma, \Pi \fCenter \Delta, \Lambda$
  \RightLabel{\CutCS}
  \BinaryInf$\Gamma, \Pi \fCenter \Delta, \Lambda$
  \DisplayProof
  \]
  وأضف $\Gamma$ إلى يسار كل متتالية أسفله، و$\Delta$ إلى يمينها. كرر ذلك
  حتى تُستبدل جميع «استدلالات» \LeftR{\lforall} التي تحتوي مقدمتها على
  $!B(t)$ معلّمة باستدلال \CutCS{} على $!B(t)$ ما. وتكون المتتالية النهائية
  لــ!!{proof} الناتج (وهو الآن !!{proof} صحيح) من الصورة
  $\Gamma, \dots, \Gamma \Sequent \Delta, \dots, \Delta$. طبّق مقبولية
  الانكماش لتحوّله إلى !!a{proof} لــ$\Gamma \Sequent \Delta$، واستبدل به
  $\pi$. وهكذا استبدلنا قطعًا واحدًا على $\lforall[x][!B(x)]$ بقطوع
  (قد تكون كثيرة) على $!B(t)$، لكنها كلها أدنى رتبة. وتبقى أي قطوع كانت
  موجودة أصلًا في $\pi_1$ أو $\pi_2$، لكنها أيضًا كلها أدنى رتبة.
\end{enumerate}
\end{proof}

\begin{prob}
تحقق من الحالة $!A \ident (!B \land !C)$ في برهان
\olref{lem:max-cut-red-G3c}. أي بيّن أنه إذا وجد !!a{proof}~$\pi$ ينتهي بــ
\[
\AxiomC{}
\RightLabel{$\pi_1$}
\Deduce$\Gamma \fCenter \Delta, !B \land !C$
\AxiomC{}
\RightLabel{$\pi_2$}
\Deduce$!B \land !C, \Gamma \fCenter \Delta$
\RightLabel{\CutCS}
\BinaryInf$\Gamma \fCenter \Delta$
\DisplayProof
\]
وكان كل من $\pi_1$ و$\pi_2$ لا يحتوي إلا قطوعًا رتبتها
$<\depth{!B \land !C}$، فهناك !!a{proof}~$\pi'$ لــ$\Gamma \Sequent
\Delta$ لا يحتوي إلا قطوعًا رتبتها $<\depth{!B \land !C}$.
\end{prob}

% \begin{prob}
% تحقق من الحالة $!A \ident \lexists[x][!B(x)]$ في برهان
% \olref{lem:inv-G3c-cut}.
% \end{prob}


تثبت \olref{lem:max-cut-red-G3c} أنه يمكننا استبدال استدلال \CutCS{} ذي
الرتبة الأعلى في !!a{proof} باستدلالات \CutCS{} أدنى رتبة. ويمكننا مرة
أخرى استعمال هذه المبرهنة لإزالة أي عدد من استدلالات \CutCS{} من
!!a{proof}، وذلك بتطبيقها تباعًا على !!{proof}s الفرعية العليا المنتهية
بــ\CutCS{} ذي الرتبة القصوى.

\begin{thm}\ollabel{thm:cut-elim-G3c-largest}
يمكن تحويل أي !!{proof}~$\pi$ في $\Log{G3c} + \CutCS$ إلى برهان في~\Log{G3c}.
\end{thm}

\begin{proof}
  نثبت أولًا أنه يمكن إزالة كل القطوع ذات الرتبة القصوى في~$\pi$. لتكن
  $d$ الرتبة القصوى للقطوع الواقعة في~$\pi$، وليكن $n$ عدد القطوع ذات
  الرتبة~$d$. البرهان بالاستقراء على~$n$. إذا كان $n=0$ فلا شيء نثبته.
  وإذا كان $n>0$، فاختر قطعًا علويًا رتبته~$d$، وليكن $\pi_1$ هو
  !!{proof} الفرعي المنتهي به. وبحسب \olref{lem:max-cut-red-G3c} يوجد
  !!a{proof}~$\pi_1'$ للمتتالية النهائية نفسها، وكل قطوعه رتبتها $<d$.
  استبدل $\pi_1$ بــ$\pi_1'$ في~$\pi$. ينتج من ذلك !!a{proof} يحتوي
  $n-1$ قطعًا من الرتبة~$d$، فينطبق فرض الاستقراء.

  والآن تتبع النتيجة بالاستقراء على~$d$. فإذا كان $d=0$، كانت جميع القطوع
  ذرية، ويمكن إزالة جميعها بحسب النتيجة السابقة. وإذا كان $d>0$، تعطينا
  النتيجة السابقة برهانًا استُبدلت فيه كل القطوع ذات الرتبة~$d$ بقطوع
  رتبتها~$<d$. ولما كانت $d$ هي الرتبة القصوى للقطوع في~$\pi$، أصبح لدينا
  برهان رتبة قطوعه القصوى $<d$، فينطبق فرض الاستقراء.
\end{proof}

\end{document}
